% a01_H0rel part 6x2: the replaced steps (b)-(d) at block factor 5 (input from a01_H0rel_p6.tex after a01_p6x1)
\subparagraph{(b) The cut-off $\zeta^{(i)}$.} In \cite[(2.44)]{Ba2} $A^{(i)}=a_i\ell_i^{-2}\zeta^{(i)}G_iQ_i^*A$ with $\supp\zeta^{(i)}(\cdot,y)\subset\{|x-y|<r_i-2M\}$, $\zeta^{(i)}(x,y)=1$ for $|x-y|\le\frac12r_i$. The printed proof of \cite[(2.60)]{Ba2}, $\partial^\eta A^{(i)}=O(p_i)$ on $\Omega$, bounds the first two terms of \cite[(2.61)]{Ba2} by the restrictions \cite[(2.55)]{Ba2} at every coarse point $y'$ with $\zeta^{(i)}(x,y')\ne0$; \cite[(2.55)]{Ba2} is available on $\Lambda_{-1}^{(i-1)\prime}$, which contains $\Omega$ at depth $\ge3r(\ell_{i-1})/L-2$ (\cref{lem:a01:X0}(a), three layers). The containment $3r(\ell_{i-1})/L-2\ge r_i-2M$ has $\alpha=1$, $\beta=3$: it holds at every scale iff $L\le3$. The gap is real: by \cite[(2.9)]{Ba2} the restrictions survive the partial resummation only on $\Lambda'_{-1}$, $\Lambda^*_{-1}$, $\Lambda_{-1}$, so at a large block which stayed large at every finer scale $A(y')$ is unrestricted and the contribution $a_i\zeta^{(i)}(x,y')(G_iQ_i^*)(x,y')A(y')$ has no bound; hence \cite[(2.60)]{Ba2} as printed is not established in a collar of width $\approx r_i-3r(\ell_{i-1})/L$ along $\partial\Omega$, while \cite[p.~571]{Ba2} uses it on all of $\Omega$ and \cite[Prop.~1]{Ba3} assumes it on \cite[p.~412]{Ba3}. The same cut-off enters the split \cite[(2.89)]{Ba2} at the next scale.

\begin{proposition}[the shrunk cut-off]\label{prop:a01:Zth}
Replace $\zeta^{(i)}$ in \cite[(2.44)]{Ba2}, at every scale, by $\zeta^{(i)}_\theta$ with $\supp\zeta^{(i)}_\theta(\cdot,y)\subset\{|x-y|<\theta r_i-2M\}$, $\zeta^{(i)}_\theta(x,y)=1$ for $|x-y|\le\frac12\theta r_i$, $|\partial^\eta\zeta_\theta|\le1$, $|\partial^\eta\partial^\eta\zeta_\theta|\le2$ (a $C^{1,1}$ radial profile), with $\theta=\theta_5\defeq\frac25$; the first-step cut-off of \cite[p.~561]{Ba2} already has range $r/(2L)$ in the units of the next scale, so this returns to the first-step convention. Let $R\ge\max\{10(M+1),\,M_\theta+2,\,(M_\theta+7)/11.6\}$, where $\theta_{i+1}$ equals $1$ on $B^i(\Lambda_2^{(i)})$ and $0$ at distance $\ge M_\theta$ from it ($M_\theta<M$ by \cite[p.~567]{Ba2}, where $\theta_{i+1}$ ``varies smoothly'' on a slice of thickness $<M$; we take $M_\theta=2M/3$ and a $C^{1,1}$ profile). Then at $L=5$:
\begin{enumerate}[label=\textup{(Z\arabic*)},nosep,leftmargin=*]
\item for $y\in\Lambda_2^{(i-1)\prime}$, $x\in B^i(y)$, the $\zeta_\theta$-ball $\{|x-y'|<\theta r_i-2M\}$ lies in $\Lambda_{-1}^{(i-1)\prime}$, where \cite[(2.55)]{Ba2} holds;
\item \cite[(2.59)--(2.60), (2.64)]{Ba2} hold with $O(p_i)$ replaced by $O_L(p_i)$, the constant containing $p_{i-1}/p_i\le(1+\ln L)^p$;
\item the terms carrying $1-\zeta^{(i)}_\theta$ in \cite[(2.61), (2.64)]{Ba2}, the terms carrying $1-\zeta^{(i)}_\theta$, $1-\zeta^{(i+1)}_\theta$ in \cite[(2.88)--(2.90)]{Ba2}, and the second term of \cite[(2.89)]{Ba2} (the one with $\delta C^{(i)}_{\Lambda_0}$) are $O(\ell_i^\kappa)$ for every $\kappa$, with their derivatives and the $\alpha_0$-H\"older quotients of their derivatives;
\item the transition width $\frac12\theta r_i-2M\ge2$ required by the $C^{1,1}$ radial profile with $|\partial^\eta\zeta_\theta|\le1$ and $|\partial^\eta\partial^\eta\zeta_\theta|\le2$ holds;
\item the split \cite[(2.89)]{Ba2} holds \emph{as printed}, with the full $B$ in its fifth term, and the $\zeta^{(l)}_\theta$-balls around $\supp\theta_l$ lie in $\Lambda_{-1}^{(l-1)\prime}$ at every scale $l$.
\end{enumerate}
No $\zeta$ obeying the printed \cite[(2.44)]{Ba2} makes \cite[(2.89)]{Ba2} exact at $L=5$ for small $\ell_i$ (the minimal printed support $\frac12r$ gives $L\theta=\frac52>2$), so the shrunk cut-off is necessary at $L=5$ if \cite[(2.89)]{Ba2} is kept in Ba{\l}aban's form.
\end{proposition}
\begin{proof}
(Z1): $\theta r_i\le3r(\ell_{i-1})/L-2+2M$ is \cref{lem:a01:X0}(b) with $\alpha=\frac25$, $\beta=3$, $\alpha L=2\le3$, margin $\ge r(\ell_{i-1})/5-2\ge0$ for $R\ge10$. (Z2): in \cite[(2.61)]{Ba2} with $\zeta_\theta$ the first term sums over $y'$ in the $\zeta_\theta$-ball, where by (Z1) $|A(y')-A(y)|\le c_1p_{i-1}L|y'-y|$ (a path of $\le L|y'-y|$ bonds with $|\partial A|\le c_1p_{i-1}$) and \cite[Prop.~2.2]{Ba2} gives the kernel bound $c_{\mathrm{Ba}}e^{-\delta_0|x-y'|}$, so the term is $O(p_{i-1})=O_L(p_i)$; the third term is the exact \cite[(2.62)--(2.63)]{Ba2} ($G_iQ_i^*1$ is constant); the second is a tail of (Z3); \cite[(2.64)]{Ba2} is the same for the derivative, with $\partial^\eta(G_iQ_i^*1)=0$. (Z3): the $(1-\zeta_\theta)$ term of \cite[(2.61)]{Ba2} multiplies the single constant $A(y)$, $|A(y)|\le c_1p_{i-1}/(\lmu\ell_{i-1})$, with a kernel supported at $|x-y'|\ge\frac12\theta r_i$ and bounded by $c_{\mathrm{Ba}}e^{-\delta_0|x-y'|}$, so it is $\le Cr_i^de^{-\delta_0\theta r_i/2}p_{i-1}/(\lmu\ell_{i-1})\le\ell_i^\kappa$ for small $\ell_i$ (\cref{lem:a01:X0}(c)); in \cite[(2.89)]{Ba2} the third and fourth terms carry $1-\zeta^{(i)}_\theta$, $1-\zeta^{(i+1)}_\theta$ acting on $\Lambda_0^{(i)\prime}B$ with $|B|\le c_1p/(\lmu\ell)$ there, and the same bound gives them, their derivatives and H\"older quotients (kernel by kernel: \cite[(2.34), (2.36)]{Ba1} for $C^{(i)}$, $\delta C_{\Lambda_0}$, the composition identity of (Z5) for $G^\eta_{i+1}Q_{i+1}^*$, the $C^{1,1}$ profile at scales $i$, $i+1$); the second term $T_2(x)=a_iaL^{-2}\theta_{i+1}(x)\sum_y\zeta^{(i)}(x,y)(G_iQ_i^*)(x,y)(\delta C^{(i)}_{\Lambda_0}Q^*\Lambda_0^{(i)\prime}B)(y)$ has, for $x\in\supp\theta_{i+1}$ and $\zeta(x,y)\ne0$, $\dist(y,\Lambda_0^{(i)c})>r_i-M_\theta-2$ (two layers of \cite[(2.8)]{Ba2} minus the range of $\zeta$), and $\delta C_{\Lambda_0}(y,y')$ decays exponentially in $|y-y'|$ and in the distances of $y$, $y'$ to $\Lambda_0^{(i)c}$ \cite[(2.35)]{Ba1}, so $|T_2|+|\partial^\eta T_2|\le Cr_i^de^{-\delta'(r_i-M_\theta-2)}p/(\lmu\ell)\le C_\kappa\ell_i^\kappa$. (Z4): $\frac12\theta r_i-2M\ge2$ iff $r_i\ge(4+4M)/\theta=10(M+1)$, and $r_i\ge R$. (Z5): on the unit lattice of scale $i$, with $Q$ the one-step average, $P_i=Q_i^*Q_i$, $P_{i+1}=Q_i^*Q^*QQ_i$, $G_i=(H+a_iP_i)^{-1}$, $G^\eta_{i+1}=(H+a_{i+1}L^{-2}P_{i+1})^{-1}$, $\Delta^{(i)}=a_i-a_i^2Q_iG_iQ_i^*$, $C^{(i)}=(aL^{-2}Q^*Q+\Delta^{(i)})^{-1}$ and $a_{i+1}=aa_i/(a_i+aL^{-2})$ \cite[(2.13), (2.21), (2.31)]{Ba1}, one has the composition identity $a_iaL^{-2}G_iQ_i^*C^{(i)}Q^*=a_{i+1}L^{-2}G^\eta_{i+1}Q^*_{i+1}$: apply $H+a_{i+1}L^{-2}P_{i+1}$ to the left side, use $HG_i=1-a_iP_iG_i$, $1-a_iQ_iG_iQ_i^*=\Delta^{(i)}/a_i$, $\Delta^{(i)}C^{(i)}=1-aL^{-2}Q^*QC^{(i)}$ and $Q^*QQ^*=Q^*$; the coefficient of $Q^*QC^{(i)}Q^*$ vanishes iff $a_{i+1}(a_i+aL^{-2})=aa_i$, which is \cite[(2.13)]{Ba1}, and the remaining coefficient is $a_{i+1}L^{-2}$ by the same relation (this is \cite[(2.41)]{Ba1} in adjoint, rescaled form). By \cite[(2.80)]{Ba2} and the Dirichlet condition \cite[(2.32)]{Ba1}, $C^{(i)}_{\Lambda_0}Q^*B=C^{(i)}_{\Lambda_0}Q^*\Lambda_0^{(i)\prime}B$; writing $C_{\Lambda_0}=C^{(i)}+\delta C_{\Lambda_0}$ and $\zeta=1-(1-\zeta)$ and using the identity, $\theta_{i+1}A^{(i)}$ equals the five terms of \cite[(2.89)]{Ba2} with $\Lambda_0^{(i)\prime}B$ in the fifth term, identically in $A'$, $B$; the printed form differs by the defect $a_{i+1}L^{-2}\theta_{i+1}\zeta^{(i+1)}_\theta G^\eta_{i+1}Q^*_{i+1}\Lambda_0^{(i)\prime c}B$, which vanishes when $\zeta^{(i+1)}_\theta(x,y)=0$ for all $x$ with $\theta_{i+1}(x)\ne0$ and $y\notin\Lambda_0^{(i)\prime}$, i.e.\ when $L(\theta r_{i+1}-2M)\le2r_i-M_\theta-2-L$ (such $x$ lie within $M_\theta+1$ of $\Lambda_2^{(i)}$, such $y$ within $L$ of $\Lambda_0^{(i)c}$, and $\dist(\Lambda_2^{(i)},\Lambda_0^{(i)c})>2r_i-1$). With $L\theta_5=2$ this reads $2(r_i-r_{i+1})\ge M_\theta+7-10M$, and $r_i-r_{i+1}=R\ln5\,(2u_{i+1}+\ln5)\ge5.81R$. The balls around $\supp\theta_l$ reach at most $2r_l-10M+M_\theta+1$ beyond $\Lambda_2^{(l-1)}$ (scale-$(l-1)$ units), while $\Lambda_{-1}^{(l-1)}$ contains the $(3r(\ell_{l-1})-2)$-neighbourhood of $\Lambda_2^{(l-1)}$ (\cref{lem:a01:X0}(a)): $2r_l+M_\theta+8\le3r(\ell_{l-1})+10M$ holds since $r_l\le r(\ell_{l-1})$ and $r(\ell_{l-1})\ge R\ge M_\theta+2$. For the printed $\theta=1$ the defect condition has $L\theta=L$, $\beta=2$, so it holds at every scale iff $L\le2$; on the reading that absorbs the defect into $\tilde B^{(i)}$, the $\zeta^{(i+1)}$-balls must stay inside $\Lambda_{-1}^{(i)\prime}$, $\beta=3$, hence $L\le3$.
\end{proof}
The other uses of \cite[Lemma~2.3]{Ba2} are covered by (Z2): \cite[(2.68)]{Ba2} on the box of (c) below; \cite[(2.95)--(2.98)]{Ba2} at each scale $l$ together with (Z5); the regularity of $\tilde B$ needed by \cite[Prop.~1]{Ba3} at the background $\theta_iA^{(i)}$ on $\Omega$, where $\theta_i=1$; and the first-step use \cite[(2.25)]{Ba2}, whose cut-off has $\theta=1/(2L)<3/L$. On the consumer side, the all-small-field term of \cite[(3.4)]{K1} is affected only through the cut background, and the difference between the actions at $A^{(j)}$ and at $A_\theta$ is $O(\eta^{s'})|X^{(j)}|$ on $\supp\chi_j$, as used in the proof of \cref{prop:a01:upper}.

\subparagraph{(c) The boxes of \cite[Lemma~2.4]{Ba2}.} The printed proof \cite[p.~572]{Ba2} takes $\square_1,\square_2$ as the sums of large blocks of $\Lambda_7^{(i-1)\prime}$ within $2r_i$, $4r_i$ of $y$ and $\square=B^i(\square_2)$. For $y\in\Lambda_7$ at distance $O(1)$ from $\Lambda_7^c$ the step \cite[(2.67)]{Ba2} is not justified: the dropped source $\Lambda_6\phi-\square_1\phi$ contains $\phi$ on $\Lambda_6\setminus\Lambda_7$ at distance $O(1)$ from $x$, where $a_iG_iQ_i^*$ is $O(1)$, so this error is $O(\Upsilon_i)$, $\Upsilon_i\defeq c_1(\lambda\ell_{i-1})^{-1/4}p_{i-1}$, not $O(\ell_i^\kappa)$; the domain change $\Omega\to\square$ happens at distance $O(1)$ from $x$, so \cite[(2.26)]{Ba1} gives no smallness; and $x$ may lie within $R_{\mathrm{Ba}}$ of $\partial\square$, where \cite[(2.24)--(2.25)]{Ba1} need a parallelepiped with a constant field near its boundary, while $\Lambda_7\cap(\text{cube})$ has re-entrant edges. An untruncated box of radius $4r_i$ inside $\Omega$ would need $4r_i\le5\rho_i$, which fails for small $\ell_i$ at every $L\ge2$ ($\alpha=4$, $\beta=5$). The only consumer is the quartic convergence factor \cite[(2.78)--(2.79)]{Ba2}, a sum over $y\in\Lambda_7^{(i-1)\prime}\cap\Lambda_7^{(i)c}$, which contains every point of $\Lambda_7$ within $2r_i$ of $\Lambda_7^c$: the edge points are exactly the ones used.

\begin{lemma}[{\cite[Lemma~2.4]{Ba2}} with a complete proof, every $L$]\label{lem:a01:lem24}
Assume \cite[(2.55)]{Ba2}, and: $5(r(\eps_0^{\mathrm{Ba}})-s_b)/L\ge R_{\mathrm{Ba}}+c_d+2s_b$, $4r(\eps_0^{\mathrm{Ba}})\ge R_{\mathrm{Ba}}+c_d$, $R\ge10s_b+10$, with $c_d\defeq\sqrt d\,s_b$; $\eps_0^{\mathrm{Ba}}$ small enough that $r(\ell_{i-1})/r(\ell_i)\le2L$ for all $i$ (at $L=5$, $\varrho=2$ this is $(1+\ln5)^2\le10$, true for all $\eps_0^{\mathrm{Ba}}\le1$); the mass condition $\Upsilon_im^2\ell_i^2(G_i(\square,0)Q_i^*1)(x)=O(p_i)$ (Ba{\l}aban's own use of \cite[(2.75)]{Ba2}); and $Ca_i\Upsilon_i(1+r_i)^{d-1}\exp(-\delta_0[(r(\ell_{i-1})-s_b)/L-c_d])\le\ell_i^\kappa$, which holds for all small $\eps_0^{\mathrm{Ba}}$ at every $R>0$ because $\varrho>1$. Then \cite[(2.65)--(2.66)]{Ba2} hold for every $y\in\Lambda_7^{(i-1)\prime}$, $x\in B^i(y)$, $b\subset B^i(\Lambda_7^{(i-1)\prime})$, with constants independent of $\eps$, $i$, the lattice and the position of $y$ relative to $\partial\Lambda_7$.
\end{lemma}
\begin{proof}
Put $s_i\defeq\min\{4r_i,5\rho_i-2s_b\}$, let $\square_2$ be the cube of large blocks of sup-radius $s_i$ around the large block of $y$, $\square\defeq B^i(\square_2)$, and $\square_1\phi\defeq1_{Q(y,s_i/2)}\Lambda_6\phi$ with $Q(y,s)$ the sup-ball. Since $\dist(\Lambda_7,\Lambda_2^c)\ge5\rho_i$ (\cref{lem:a01:X0}(a)), $\square\subset\Omega$; $\square$ is a rectangular parallelepiped and a union of large blocks, and $s_i-c_d\ge R_{\mathrm{Ba}}$ at every $i$ by the hypotheses. (1) \emph{Source cut:} $\Lambda_6\phi-\square_1\phi$ is supported at distance $\ge s_i/2-c_d$ from $x$, and $\dist(x,\Omega^c)\ge5\rho_i-c_d\ge R_{\mathrm{Ba}}$, so \cite[(2.25)]{Ba1} on $\Omega$ with the sup bound $\Upsilon_i$ on $\phi$ gives $|a_i(G_i(\Omega,A^{(i)})Q_i^*(\Lambda_6\phi-\square_1\phi))(x)|\le Ca_ie^{-\delta_0(s_i/2-c_d)}\Upsilon_i$. (2) \emph{Domain change:} $G_i(\square,A)=G_i(\Omega,A)+\delta G_i(\square,\Omega,A)$, and by \cite[(2.26)]{Ba1} the second term applied to $Q_i^*\square_1\phi$ at $x$ carries $e^{-\delta_0(3s_i/2-2c_d)}$; the same for $D^\eta$ by \cite[(2.58)]{Ba2}. Hence $\phi^{(i)}(x)=a_i(G_i(\square,A^{(i)})Q_i^*(A^{(i)})\square_1\phi)(x)+O(a_i\Upsilon_ie^{-\delta_0(s_i/2-c_d)})$, and likewise for the covariant derivative: this is \cite[(2.67)]{Ba2} with the boxes chosen above. (3) \emph{Constant background:} $A_0\defeq A^{(i)}(y)$; joining points of the convex box by straight lattice paths and using \cite[(2.60)]{Ba2} (\cref{prop:a01:Zth}(Z2)) on $B^i(\Lambda_2^{(i-1)\prime})\supset\square$, $|A^{(i)}(z)-A_0|\le O(p)\sqrt d\,s_i\le O(pr_i)$ on $\square$; the expansion \cite[(3.44)]{Ba1} and \cite[Prop.~2.2]{Ba1} on the parallelepiped $\square$ replace $A^{(i)}$ by $A_0$ with error $O(\ell_i^{\kappa_0})$, as in \cite[(2.68)]{Ba2}. (4) \emph{Gauge-out:} $A_0=\partial^\eta\lambda_0$ with $\lambda_0(z)=A_0\cdot(z-y)$ on $\square$, so $U(A_0(\Gamma_{z,y}))=e^{ieA_0\cdot(z-y)}$ for every contour inside $\square$, $G_i(\square,\partial\lambda_0)=U(\lambda_0)G_i(\square,0)U(\lambda_0)^*$ and $Q_i(\partial\lambda_0)=U(\lambda_0)_{\mathrm{coarse}}Q_i(0)U(\lambda_0)^*$: \cite[(2.69)--(2.74)]{Ba2}; legitimate because $\square$, of diameter $\le8r_i+s_b$, does not wrap around (on the cube, no condition). Write $\phi'\defeq U(\lambda_0)^{-1}\phi$. (5) \emph{Reproduction of constants:} by \cite[(2.75)]{Ba2}, $a_iG_i(\square,0)Q_i^*1=1-m^2\ell_i^2G_i(\square,0)Q_i^*1$, hence $a_i(G_i(\square,0)Q_i^*\square_1\phi')(x)-\phi'(y)=a_i(G_i(\square,0)Q_i^*[\square_1(\phi'-\phi'(y))])(x)-\phi'(y)a_i(G_i(\square,0)Q_i^*1_{\square\setminus\square_1})(x)-\phi'(y)m^2\ell_i^2(G_i(\square,0)Q_i^*1)(x)$. The first term is $O(p)$: on the convex cube $Q(y,s_i/2)\subset\Lambda_{-1}$ the restriction $|D\phi|\le c_1p$ gives $|\phi'(z)-\phi'(y)|\le O(p)|z-y|$ and the kernel decays exponentially. The second: $\square\setminus\square_1\subset(Q(y,s_i)\setminus Q(y,s_i/2))\cup(Q(y,s_i/2)\setminus\Lambda_6)$ lies at distance $\ge\min\{s_i/2,\rho_i\}-c_d$ from $x$, so it is $\le c_{\mathrm{Ba}}e^{-\delta_0(\min\{s_i/2,\rho_i\}-c_d)}\Upsilon_i$ by \cite[(2.25)]{Ba1} on the parallelepiped. The third is the mass-defect term of the printed proof, $O(p_i)$ under the mass condition. (6) \emph{Derivatives:} a bond $b\subset B^i(y)$ lies at distance $\ge s_i-c_d\ge R_{\mathrm{Ba}}$ inside $\square$ and $\ge5\rho_i-c_d$ inside $\Omega$, so the printed derivative bounds apply at $b$ and (1)--(5) give \cite[(2.66)]{Ba2} in the form \cite[(2.77)]{Ba2}; no re-entrant edge is approached. \emph{Errors:} $s_i/2\ge\rho_i$ iff $\rho_i\le2r_i$ and $3\rho_i\ge2s_b$, the first by $r(\ell_{i-1})/r_i\le2L$, the second by $R\ge10s_b+10$; so every error is $\le Ca_i\Upsilon_i(1+r_i)^{d-1}e^{-\delta_0(\rho_i-c_d)}$, which is $\le\ell_i^\kappa$ by the last hypothesis; the rest of the printed proof, \cite[(2.69)--(2.77)]{Ba2}, runs on $\square$. Nothing depends on $L$ except through the named constants.
\end{proof}
At $L=5$ the hypotheses read $r(\eps_0^{\mathrm{Ba}})\ge R_{\mathrm{Ba}}+(3+\sqrt3)s_b$, $r(\ell_{i-1})/r_i\le6.81\le10$, $R\ge10s_b+10$: lower bounds on $R$ and caps on $\eps_0^{\mathrm{Ba}}$, none restricting $L$. On the cube the same device (a box of large blocks inside the propagator domain about the collar cubes, with the cap inherited from the previous step) is the King box of D8.
